% ============================================================================
%  ECE 431 -- Sheet 2 Solutions
%  Style: sheet-answering-style.md
% ============================================================================

\begin{question}[Question 1]Consider a HeNe laser source of area $A$ that emits 1\,mW of red-light
($h\nu\approx 2.0$\,eV) into a solid angle $\lambda^{2}/A$. Calculate
\begin{enumerate}[label=(\alph*)]
  \item The number of photons emitted per second by this laser source.
  \item The solid angle of emission if $A=80$\,mm$^{2}$.
\end{enumerate}
\end{question}

\textbf{\large Solution --- Question 1}

\textbf{Part (a): Photon emission rate.}

Each photon carries energy $h\nu$, so the power is the rate times the photon energy. Rearranging:

\begin{equation*}
  P = \dot{N}\cdot h\nu
  \quad\Longrightarrow\quad
  \dot{N} = \frac{P}{h\nu} = \frac{1\times10^{-3}}{2.0\times1.602\times10^{-19}}
\end{equation*}

\begin{equation*}
  \boxed{\dot{N} = 3.12\times10^{15}\;\text{photons/s.}}
\end{equation*}

\textbf{Part (b): Solid angle with $A = 80\,\text{mm}^{2}$.}

First recover the wavelength from the photon energy $h\nu = 2.0\,\text{eV}$:

\begin{equation*}
  \lambda = \frac{hc}{h\nu} = \frac{(6.626\times10^{-34})(3\times10^{8})}{2.0\times1.602\times10^{-19}} = 620\,\text{nm}.
\end{equation*}

Then the emission solid angle is:

\begin{equation*}
  \Omega = \frac{\lambda^{2}}{A} = \frac{(620\times10^{-9})^{2}}{80\times10^{-6}}
\end{equation*}

\begin{equation*}
  \boxed{\Omega = 4.805\times10^{-9}\;\text{sr.}}
\end{equation*}

This extremely small solid angle is the hallmark of laser directionality: the diffraction limit $\lambda^2/A$ is orders of magnitude smaller than the $2\pi$\,sr of a typical LED.

% ============================================================================

\begin{question}[Question 2]Starting from the spectral energy density (J/m$^{3}$-Hz) of the black
body, find an expression of the spectral intensity (watts/m$^{2}$-Hz) emitted by that body.
Then find the photon flux density $\Phi$ (no.\ of photons/m$^{2}$-sec) emitted within
a bandwidth $\Delta\nu$ at frequency $\nu$. Then calculate:

\begin{enumerate}[label=(\alph*)]
  \item The number of photons of red light emitted per second by a black body (thermal light
        source) of temperature $T=1000$\,K within a bandwidth of 10\,kHz and a solid angle
        $\lambda^{2}/A$, assuming that it has the same area as the laser source in Problem (1).
  \item The temperature of the black body required to emit the same number of photons per
        second as in Problem (1).
\end{enumerate}
\end{question}

\textbf{\large Solution --- Question 2}

\textbf{Spectral energy density and intensity.}

The Planck spectral energy density (energy per unit volume per unit frequency) is:

\begin{equation*}
  \xi_\text{BB}(\nu) = \frac{8\pi h}{\lambda^{3}}\cdot\frac{1}{e^{h\nu/k_BT}-1}.
\end{equation*}

Each photon travels at speed $c$, so the spectral intensity (power flowing through unit area per unit frequency) is obtained by multiplying by $c$:

\begin{equation*}
  \rho(\nu) = \xi_\text{BB}(\nu)\cdot c = \frac{8\pi hc}{\lambda^{3}}\cdot\frac{1}{e^{h\nu/k_BT}-1}.
\end{equation*}

Dividing by the photon energy $h\nu$ and multiplying by $\Delta\nu$ gives the photon flux density (photons per m$^{2}$ per second in bandwidth $\Delta\nu$). Using $c/(\lambda^{3}\nu) = 1/\lambda^{2}$:

\begin{equation*}
  \boxed{\Phi_\nu = \frac{8\pi}{\lambda^{2}}\cdot\frac{\Delta\nu}{e^{h\nu/k_BT}-1}.}
\end{equation*}

\textbf{Part (a): Photon rate from black body at $T = 1000$\,K.}

The number of photons per second collected from an emitting area $A$ into solid angle $\Omega = \lambda^{2}/A$ in bandwidth $\Delta\nu = 10^{4}$\,Hz is $\dot{N} = A\cdot\Phi_\nu = (8\pi A/\lambda^{2})\cdot\Delta\nu/(e^{h\nu/k_BT}-1) = (8\pi/\Omega)\cdot\Delta\nu/(e^{h\nu/k_BT}-1)$. With $\Omega = 4.805\times10^{-9}$\,sr and $h\nu/k_BT = 2.0\times1.602\times10^{-19}/(1.381\times10^{-23}\times1000) = 23.2$:

\begin{equation*}
  \dot{N} = \frac{8\pi}{4.805\times10^{-9}}\times\frac{10^{4}}{e^{23.2}-1}
\end{equation*}

\begin{equation*}
  \boxed{\dot{N} \approx 4355\;\text{photons/s.}}
\end{equation*}

This is eleven orders of magnitude smaller than the laser rate of $3.12\times10^{15}$\,photons/s, illustrating the enormous brightness advantage of coherent laser emission over thermal radiation.

\textbf{Part (b): Temperature to match the laser.}

Setting the same formula equal to $3.12\times10^{15}$\,photons/s and solving for $T$:

\begin{equation*}
  3.12\times10^{15} = \frac{8\pi}{4.805\times10^{-9}}\times\frac{10^{4}}{e^{2.0\,\text{eV}/k_BT}-1}
  \quad\Longrightarrow\quad
  e^{2.0\,\text{eV}/k_BT} - 1 \approx 1.67\times10^{-6}.
\end{equation*}

For such a small argument, $e^x - 1 \approx x$, so $2.0\,\text{eV}/k_BT \approx 1.67\times10^{-6}$, giving:

\begin{equation*}
  \boxed{T = \frac{2.0\,\text{eV}}{k_B\times1.67\times10^{-6}} \approx 1.396\times10^{6}\;\text{K.}}
\end{equation*}

No real material survives at such a temperature, which is why laser sources cannot be replicated by thermal means.

% ============================================================================

\begin{question}[Question 3]Show that the Einstein model of light-matter interaction applies for
broad-band light sources without the need to include the line-shape function. When must the
line-shape function be incorporated?
\end{question}

\textbf{\large Solution --- Question 3}

The Einstein transition rates involve an integral of the spectral energy density $\xi(\nu)$ weighted by the atomic line-shape function $g(\nu)$. For spontaneous emission (independent of the field) the rate per atom is simply $1/\tau_\text{spon}$. For absorption and stimulated emission:

\begin{equation*}
  \left.\frac{dN_2}{dt}\right|_\text{absorp}
  = \int_\text{all}\frac{\lambda^{3}N_1}{8\pi h\tau_\text{spon}}\,\xi(\nu)\,g(\nu)\,d\nu.
\end{equation*}

For a \emph{broadband} source, $\xi(\nu)$ varies slowly over the narrow atomic linewidth $\Delta\nu$. It may therefore be factored outside the integral:

\begin{equation*}
  \left.\frac{dN_2}{dt}\right|_\text{absorp}
  = \frac{\lambda^{3}N_1}{8\pi h\tau_\text{spon}}\,\xi(\nu)
    \underbrace{\int_\text{all}g(\nu)\,d\nu}_{=\;1}
  = \frac{\lambda^{3}N_1}{8\pi h\tau_\text{spon}}\,\xi(\nu).
\end{equation*}

Because $g(\nu)$ is normalised to unity, it drops out entirely:

\begin{equation*}
  \boxed{\left.\frac{dN_2}{dt}\right|_\text{absorp} = \frac{\lambda^{3}N_1}{8\pi h\tau_\text{spon}}\,\xi(\nu),
  \qquad
  \left.\frac{dN_2}{dt}\right|_\text{stim} = -\frac{\lambda^{3}N_2}{8\pi h\tau_\text{spon}}\,\xi(\nu).}
\end{equation*}

The same argument applies to stimulated emission. The line-shape function must be \emph{retained} for \textbf{narrowband (monochromatic) sources}, where $\xi(\nu)$ varies sharply over the linewidth and cannot be treated as constant inside the integral.

% ============================================================================

\begin{question}[Question 4]Show that the line-shape function in the case of inhomogeneous Doppler
broadening may be written in the form
\begin{equation*}
  \bar{g}(\nu)=\int_{-\infty}^{+\infty}g_{\nu_{o}}\!\left(\nu-\nu_{o}\frac{V}{c}\right)P(V)\,dV
\end{equation*}
where
\begin{equation*}
  P(V)=\frac{1}{\sqrt{2\pi\sigma_{V}^{2}}}\,e^{-V^{2}/2\sigma_{V}^{2}}.
\end{equation*}
Then find an expression of $\bar{g}(\nu)$ in the two limits
\begin{enumerate}[label=(\arabic*)]
  \item $\sigma_{V}\to 0$ and
  \item $\Delta\nu\to 0$.
\end{enumerate}
\emph{Hint:} $\lim_{\sigma_{V}\to 0}P(V)=\delta(V)$ and
$\lim_{\epsilon\to 0}\epsilon/(\epsilon^{2}+x^{2})=\pi\delta(x)$.
\end{question}

\textbf{\large Solution --- Question 4}

\textbf{Deriving the averaged line-shape.}

A photon at frequency $\nu'$ (lab frame) is seen by an atom moving at velocity $V$ (along the propagation direction) at frequency $\nu = \nu'(1 - V/c)$. The atom resonates when this Doppler-shifted frequency matches its resonance, i.e.\ its effective resonance in the lab frame is shifted to:

\begin{equation*}
  \nu'_o = \frac{\nu_o}{1 - V/c} \approx \nu_o\!\left(1 + \frac{V}{c}\right).
\end{equation*}

Each velocity class has a Lorentzian line-shape centred on $\nu'_o$:

\begin{equation*}
  g_{\nu'_o}(\nu) = g_{\nu_o}\!\left(\nu - \nu_o\frac{V}{c}\right).
\end{equation*}

Averaging over the Maxwell--Boltzmann velocity distribution $P(V)$:

\begin{equation*}
  \boxed{\bar{g}(\nu) = \int_{-\infty}^{+\infty} g_{\nu_o}\!\left(\nu - \nu_o\frac{V}{c}\right) P(V)\,dV.}
\end{equation*}

\textbf{Limit (1): $\sigma_V \to 0$.}

As $\sigma_V \to 0$, $P(V) \to \delta(V)$. The integral collapses to a single velocity class at $V = 0$:

\begin{equation*}
  \bar{g}(\nu) = \int_{-\infty}^{+\infty} g_{\nu_o}\!\left(\nu - \nu_o\frac{V}{c}\right)\delta(V)\,dV = g_{\nu_o}(\nu).
\end{equation*}

When there is no thermal motion, the inhomogeneous average reduces to the single-atom homogeneous lineshape --- as expected.

\textbf{Limit (2): $\Delta\nu \to 0$.}

Substituting the Lorentzian $g_{\nu_o}\!\left(\nu - \nu_o V/c\right) = \frac{\Delta\nu/2\pi}{(\Delta\nu/2)^2 + (\nu - \nu_o - \nu_oV/c)^2}$ and changing variable $\nu_D = \nu - \nu_o V/c$, so $V = (c/\nu_o)(\nu - \nu_D)$ and the Gaussian becomes one over $\sigma_D = (\nu_o/c)\sigma_V$. As $\Delta\nu \to 0$ the Lorentzian becomes $\pi\delta(\nu_D - \nu)$ (by the hint), and the convolution integral yields a pure Gaussian:

\begin{equation*}
  \boxed{\bar{g}(\nu) = \frac{1}{\sqrt{2\pi}\,\sigma_D}\,
    e^{-(\nu-\nu_o)^{2}/2\sigma_D^{2}}, \qquad \sigma_D = \frac{\nu_o}{c}\sigma_V.}
\end{equation*}

In this limit of vanishing homogeneous linewidth the measured line-shape is purely Gaussian, determined entirely by the velocity spread of the ensemble.

% ============================================================================

\begin{question}[Question 5]The Doppler line-shape function of an inhomogeneously broadened laser is
given by
\begin{equation*}
  \bar{g}(\nu)=\frac{1}{\sqrt{2\pi}\,\sigma_{D}}
  \exp\!\left[-\frac{(\nu-\nu_{o})^{2}}{2\sigma_{D}^{2}}\right],
  \qquad
  \sigma_{D}=\frac{1}{\lambda_{o}}\sqrt{\frac{k_{B}T}{M}}.
\end{equation*}
\begin{enumerate}[label=(\alph*)]
  \item Show that the full-width half maximum (FWHM) of the Doppler linewidth is given by
        $\Delta\nu=\sqrt{8\ln 2}\,\sigma_{D}$.
  \item Consider the $0.6328\,\mu$m transition of Ne in He-Ne laser. If $M=20u$ where
        $u=1.66\times10^{-27}$\,kg is the atomic mass unit, then find the FWHM of the Doppler
        linewidth at $T=300$\,K.
\end{enumerate}
\end{question}

\textbf{\large Solution --- Question 5}

\textbf{Part (a): FWHM of the Doppler lineshape.}

The Gaussian lineshape peaks at $\nu = \nu_o$ with maximum value $\bar{g}_{\mathrm{max}} = 1/(\sqrt{2\pi}\,\sigma_D)$. The FWHM is the width between the two points where $\bar{g} = \bar{g}_{\mathrm{max}}/2$. Setting:

\begin{equation*}
  \frac{1}{\sqrt{2\pi}\,\sigma_D}\,e^{-(\Delta\nu/2)^{2}/(2\sigma_D^{2})} = \frac{1}{2}\cdot\frac{1}{\sqrt{2\pi}\,\sigma_D}
  \quad\Longrightarrow\quad
  e^{-(\Delta\nu/2)^{2}/(2\sigma_D^{2})} = \frac{1}{2}.
\end{equation*}

Taking the natural log and solving:

\begin{equation*}
  \frac{(\Delta\nu/2)^{2}}{2\sigma_D^{2}} = \ln 2
  \quad\Longrightarrow\quad
  \Delta\nu = 2\sqrt{2\ln 2}\,\sigma_D
\end{equation*}

\begin{equation*}
  \boxed{\Delta\nu = \sqrt{8\ln 2}\,\sigma_D.}
\end{equation*}

\textbf{Part (b): He-Ne laser at $T = 300$\,K.}

For the $\lambda_o = 0.6328\,\mu$m transition of Ne with $M = 20\times1.66\times10^{-27} = 3.32\times10^{-26}$\,kg:

\begin{equation*}
  \sigma_D = \frac{1}{\lambda_o}\sqrt{\frac{k_BT}{M}}
  = \frac{1}{0.6328\times10^{-6}}\sqrt{\frac{1.381\times10^{-23}\times300}{3.32\times10^{-26}}}
  = 5.58\times10^{8}\;\text{Hz}.
\end{equation*}

\begin{equation*}
  \boxed{\Delta\nu = \sqrt{8\ln 2}\times5.58\times10^{8} \approx 1.314\;\text{GHz}.}
\end{equation*}

This GHz-scale linewidth is much wider than the natural (Lorentzian) linewidth, confirming that room-temperature He-Ne is dominated by inhomogeneous Doppler broadening.

% ============================================================================

\begin{question}[Question 6]Show that the change in the energy of emitted photons in a transition
between two energy levels, $E_{1}$ and $E_{2}$, is given by
\begin{equation*}
  \Delta E_{\mathrm{ph}}=\frac{h(1/\tau_{1}+1/\tau_{2})}{2\pi},
\end{equation*}
where $\tau_{1}$ and $\tau_{2}$ are the lifetimes of the two energy levels. Explain how this
change in photon energy conforms with the uncertainty principle. What would $\Delta E_{\mathrm{ph}}$
be if the lower energy level is the ground state? Find the corresponding absorption cross section
at the central (resonance) frequency $\nu_{o}\equiv(E_{2}-E_{1})/h$. Under which conditions is
this absorption cross section maximum? Find an expression of this maximum.
\end{question}

\textbf{\large Solution --- Question 6}

\textbf{Photon energy spread from energy-time uncertainty.}

The energy-time form of the uncertainty principle states $\Delta E\cdot\tau \geq h/(2\pi)$. Each energy level $E_i$ has a finite lifetime $\tau_i$, giving it an energy uncertainty:

\begin{equation*}
  \Delta E_2 = \frac{h}{2\pi\tau_2}, \qquad \Delta E_1 = \frac{h}{2\pi\tau_1}.
\end{equation*}

The emitted photon energy spans from $E_{\mathrm{ph,min}} = (E_2 - \Delta E_2/2)-(E_1+\Delta E_1/2)$ to $E_{\mathrm{ph,max}} = (E_2+\Delta E_2/2)-(E_1-\Delta E_1/2)$, so the spread is:

\begin{equation*}
  \boxed{\Delta E_{\mathrm{ph}} = \Delta E_2 + \Delta E_1 = \frac{h}{2\pi}\!\left(\frac{1}{\tau_1}+\frac{1}{\tau_2}\right).}
\end{equation*}

This is precisely the natural linewidth: the photon ``inherits'' the energy uncertainty of both levels involved in the transition, which is itself a statement of the energy-time uncertainty principle applied to the transition.

\textbf{Ground-state lower level ($\tau_1 \to \infty$).}

The ground state has infinite lifetime, $\Delta E_1 = 0$, so:

\begin{equation*}
  \Delta E_{\mathrm{ph}} = \frac{h}{2\pi\tau_2} \quad\Longrightarrow\quad \Delta\nu = \frac{1}{2\pi\tau_2}.
\end{equation*}

\textbf{Absorption cross section and its maximum.}

At the resonance frequency $\nu_o = (E_2-E_1)/h$, the cross section is:

\begin{equation*}
  \sigma(\nu_o) = \frac{\lambda^{2}}{8\pi\tau}\,g_{\nu_o}(\nu_o).
\end{equation*}

For a Lorentzian with FWHM $\Delta\nu$, the peak value is $g_{\nu_o}(\nu_o) = 2/(\pi\Delta\nu)$, giving:

\begin{equation*}
  \sigma(\nu_o) = \frac{\lambda^{2}}{8\pi\tau}\cdot\frac{2}{\pi\Delta\nu} = \frac{\lambda^{2}}{4\pi^{2}\tau\Delta\nu}.
\end{equation*}

The cross section is maximised when the linewidth equals the natural linewidth $\Delta\nu = 1/(2\pi\tau)$:

\begin{equation*}
  \boxed{\sigma_{\mathrm{max}}(\nu_o) = \frac{\lambda^{2}}{4\pi^{2}\tau}\cdot 2\pi\tau = \frac{\lambda^{2}}{2\pi}.}
\end{equation*}

This ultimate limit is achieved only when no other broadening mechanism (Doppler, collisional) competes with the natural linewidth --- a condition met in cold, isolated atoms.

% ============================================================================

\begin{question}[Question 7]Consider a medium with the following characteristics: radiation time
constant $\tau_{r}=230\,\mu$s, resonance wavelength $\lambda_{o}=1.06\,\mu$m, refractive index
$n=1.82$, and a FWHM $\Delta\nu=3\times10^{12}$\,Hz of the Lorentzian line-shape function.

\begin{enumerate}[label=(\alph*)]
  \item Calculate the absorption cross section at the resonance wavelength.
  \item Calculate the atomic density required to give an attenuation constant of 1.0\,m$^{-1}$.
\end{enumerate}
\end{question}

\textbf{\large Solution --- Question 7}

\textbf{Part (a): Absorption cross section at resonance.}

For a Lorentzian lineshape the cross section at the peak is $\sigma(\nu_o) = (\lambda^2/8\pi\tau_r)\,g_{\nu_o}(\nu_o)$ where $\lambda = \lambda_o/n$ is the in-medium wavelength and $g_{\nu_o}(\nu_o) = 2/(\pi\Delta\nu)$:

\begin{equation*}
  \sigma(\nu_o) = \frac{3\lambda_o^{2}}{8\pi n^{2}\tau_r}\cdot\frac{2}{\pi\Delta\nu}
  = \frac{3\lambda_o^{2}}{4\pi^{2}n^{2}\tau_r\Delta\nu}.
\end{equation*}

Substituting $\lambda_o = 1.06\times10^{-6}$\,m, $n = 1.82$, $\tau_r = 230\times10^{-6}$\,s, $\Delta\nu = 3\times10^{12}$\,Hz:

\begin{equation*}
  \sigma(\nu_o) = \frac{3\times(1.06\times10^{-6})^{2}}{4\pi^{2}\times(1.82)^{2}\times230\times10^{-6}\times3\times10^{12}}
\end{equation*}

\begin{equation*}
  \boxed{\sigma(\nu_o) \approx 37.35\times10^{-24}\;\text{m}^{2} = 37.35\;\text{pm}^{2}.}
\end{equation*}

\textbf{Part (b): Required atomic density.}

The Beer--Lambert attenuation constant relates to the cross section and atomic density by $\alpha = \sigma N_a$. Inverting:

\begin{equation*}
  N_a = \frac{\alpha}{\sigma} = \frac{1.0}{37.35\times10^{-24}}
\end{equation*}

\begin{equation*}
  \boxed{N_a \approx 2.67\times10^{22}\;\text{m}^{-3}.}
\end{equation*}

For comparison, atmospheric gas densities are around $2.7\times10^{25}$\,m$^{-3}$, so only about 0.1\% of atoms need to be in this particular transition to produce 1\,m$^{-1}$ attenuation.

% ============================================================================

\begin{question}[Formative Question]Consider a 3D cubic optical cavity whose walls are made
of perfectly conducting metal.
\begin{enumerate}[label=(\alph*)]
  \item Show that the density of states of photons in this cavity is given by
        $M(\nu)=8\pi\nu^{2}/c^{3}$, where $\nu$ is the photon frequency and $c$ is the velocity
        of light in the medium of the cavity.
  \item If the average number of photons per state is given by the Bose-Einstein distribution
        $f(\nu)=1/(e^{h\nu/K_{B}T}-1)$, then find the average energy of photon radiation per
        unit volume per unit energy.
  \item Repeat (a) and (b) in the case of a 1D optical cavity assuming the same boundary
        conditions.
\end{enumerate}
\end{question}

\textbf{\large Solution --- Formative Question}

\textbf{Part (a): 3D density of states.}

In a cubic cavity of side $L$, the conducting-wall boundary conditions enforce standing waves with integer half-wavelengths fitting along each axis: $k_x = n_x\pi/L$, $k_y = n_y\pi/L$, $k_z = n_z\pi/L$ with $n_x,n_y,n_z = 1,2,3,\ldots$\,. The dispersion relation $k = 2\pi\nu/c$ links the mode index magnitude $|\mathbf{n}| = kL/\pi = 2\nu L/c$ to the frequency. Counting all $(n_x,n_y,n_z)$ triples in the positive octant, and multiplying by 2 for the two electromagnetic polarisations:

\begin{equation*}
  N(\nu) = 2\times\frac{1}{8}\times\frac{4\pi}{3}\left(\frac{2\nu L}{c}\right)^{\!3}
  = \frac{8\pi\nu^{3}L^{3}}{3c^{3}}.
\end{equation*}

Differentiating with respect to $\nu$ and dividing by the volume $L^3$:

\begin{equation*}
  \boxed{M(\nu) = \frac{1}{L^3}\frac{dN}{d\nu} = \frac{8\pi\nu^{2}}{c^{3}}.}
\end{equation*}

\textbf{Part (b): Average energy density (3D).}

Each mode holds on average $f(\nu) = 1/(e^{h\nu/k_BT}-1)$ photons of energy $h\nu$. The average energy per unit volume per unit frequency is:

\begin{equation*}
  \boxed{\xi(\nu) = M(\nu)\cdot h\nu\cdot f(\nu) = \frac{8\pi h\nu^{3}}{c^{3}}\cdot\frac{1}{e^{h\nu/k_BT}-1}.}
\end{equation*}

This is precisely the Planck blackbody spectral energy density, confirming that the mode-counting approach reproduces classical radiation theory.

\textbf{Part (c): 1D cavity.}

For a 1D cavity of length $L$, modes satisfy $k_n = n\pi/L$ with $n = 1,2,\ldots$\,. Since $k = 2\pi\nu/c$, the mode index is $n_\nu = 2\nu L/c$. Including two polarisations, the number of modes with frequency $<\nu$ is $N_{1\mathrm{D}}(\nu) = 2\times 2\nu L/c = 4\nu L/c$.

\begin{equation*}
  \boxed{M_{1\mathrm{D}}(\nu) = \frac{1}{L}\frac{dN_{1\mathrm{D}}}{d\nu} = \frac{4}{c}.}
\end{equation*}

The mode density is uniform in frequency --- a flat spectrum, unlike the $\nu^2$ growth in 3D. The corresponding energy density per unit length per unit frequency is:

\begin{equation*}
  \xi_{1\mathrm{D}}(\nu) = M_{1\mathrm{D}}(\nu)\cdot h\nu\cdot f(\nu)
  = \frac{4h\nu}{c}\cdot\frac{1}{e^{h\nu/k_BT}-1}.
\end{equation*}
